В ходе настоящей работы были получены следующие результаты:
\begin{itemize}
  \item Показано, что облака в пространстве Громова--Хаусдорфа
    являются собственными классами в системе аксиом фон
    Неймана--Бернайса--Геделя (NBG), а не множествами.
  \item Установлено, что расстояние Громова--Хаусдорфа между облаками
    с нетривиальным пересечением стационарных групп может принимать
    только значения \( 0 \) или \( \infty \)
  \item Доказана теорема об образе центра: при соответствии с
    конечным искажением  между облаком одноточечных пространств и
    облаком с нетривиальной стационарной группой образ центра лежит
    на расстоянии, не превосходящем \( \infty \).
  \item Исследовано облако вещественной прямой: вычислены расстояния
    до целочисленной решётки и других подпространств.
  \item Введены понятия угла и равнобедренного угла облака как новых
    инвариантов геометрии облаков.
  \item Доказано, что если облака имеют нетривиальное пересечение
    стационарных групп и различные углы, то расстояние между ними
    равно бесконечности.
\end{itemize}
\section{Заключение.}
В работе исследовано расстояние Громова--Хаусдорфа между облаками ---
классами эквивалентности метрических пространств по отношению
конечности расстояния. Рассмотрены свойства облаков с точки зрения
теории множеств, показано, что они представляют собой собственные
классы в системе аксиом фон Неймана--Бернайса--Гделя.

Изучена структура облаков с нетривиальными стационарными группами.
Установлено, что
расстояние между облаками, стационарные группы которых имеют
нетривиальное пересечение, дискретно и принимает значения \( 0 \) или
\( \infty \).

Введён и исследован новый инвариант --- угол облака, обобщающий
геометрические свойства пространства ограниченных метрических
пространств. Вычислены значения углов для облака одноточечного
пространства и облака вещественной прямой. Доказана теорема,
связывающая углы облаков с расстоянием между ними: показано, что
различие углов при нетривиальном пересечении стационарных групп
влечёт бесконечность расстояния Громова--Хаусдорфа.

Полученные результаты уточняют структуру пространства всех
метрических пространств с точки зрения метрики Громова--Хаусдорфа на
неограниченных объектах. Выявлены ограничения на возможность
конечного расстояния между облаками различной геометрической природы,
показаны новые способы классификации облаков через их угловые инварианты.
